\section{Two separated point charges}
\label{sec:c3-separated-hall}

Two point modules with different first coordinates have no extensions
between them. Their Hall correspondence can therefore be computed
before the collision terms of the affine Yangian enter.
For example, take $p=(0,0,0)$ and $q=(1,0,0)$ in $\mathbb C^3$.
The module $\mathcal O_p\oplus\mathcal O_q$ has exactly two submodules
of length one. Ordering these two submodules produces the two terms
of the Hall product. The coefficient sheaf still requires a calculation:
the potential is quadratic transverse to this family of commuting triples.

Work over $\mathbb C$, with rational constructible coefficients on
analytic quotient stacks. Cohomology means ordinary equivariant
hypercohomology, defined through the Borel construction.
There is no torus acting on the three spatial coordinates in this
calculation. The general linear groups record changes of basis of modules.
Set
\[
 \mathfrak X_n=[\operatorname{Mat}_n(\mathbb C)^3/GL_n],\qquad
 f_n(A,B,C)=\operatorname{Tr}(A[B,C]),\qquad
 \mathfrak M_n=\operatorname{Crit}(f_n)^{\mathrm{cl}}.
\]
The superscript means the classical critical substack, defined by
the three commutator equations. The coefficient complex below retains
the potential on its smooth ambient stack.
The dimension of the smooth ambient stack $\mathfrak X_n$ is $2n^2$.
Use $\phi_f^p=\phi_f[-1]$, with $\phi_0^p\mathbb Q=\mathbb Q$,
and the ambient critical-chart orientation. Write
\[
 \mathcal P_n=
 \left.\phi_{f_n}^p\mathbb Q_{\mathfrak X_n}[2n^2]
 \right|_{\mathfrak M_n}.
\]
These are rational coefficient complexes. Tate twists are forgotten.
Thus no assertion about mixed Hodge structures depends on an implicit
choice of a square root of the Tate object.
The perverse shift and the ambient quotient are both part of this
definition. In particular, $\mathcal P_1=\mathbb Q_{\mathfrak M_1}[2]$,
where $\mathfrak M_1=\mathbb C^3\times B\mathbb G_m$.

Let
\[
 D=\{(p,q)\in\mathbb C^3\times\mathbb C^3:p_x\ne q_x\},\qquad
 \mathfrak D=D\times B(\mathbb G_m^2).
\]
The involution $\tau$ exchanges $p,q$ and the two line bundles.
Let $\mathfrak S\subset\mathfrak M_2$ be the open substack where $A$
has distinct eigenvalues. The direct-sum morphism is
\[
 s:\mathfrak D\longrightarrow\mathfrak S,\qquad
 (\mathcal O_p\otimes L_1,\mathcal O_q\otimes L_2)
 \longmapsto(\mathcal O_p\otimes L_1)\oplus(\mathcal O_q\otimes L_2).
\]

\begin{lemma}[The extension correspondence on the separated locus]
\label{lem:c3-separated-correspondence}
The morphism $s$ identifies $\mathfrak S$ with
$[\mathfrak D/S_2]$ and is representable finite \emph{\'{e}tale} of
degree two. The stack $\mathfrak E$ of length-one submodules in
objects of $\mathfrak S$ fits into a correspondence
\[
 \mathfrak D\xleftarrow{\ a\ }\mathfrak E
 \xrightarrow{\ b\ }\mathfrak S
\]
in which $a$ is an isomorphism and $b=s\circ a$.
Here $a$ records the submodule first and the quotient second.
\end{lemma}

\begin{proof}
Differentiation of $f_2$ in $A,B,C$ gives $[B,C],[C,A],[A,B]$.
The trace pairing on matrices is nondegenerate. Hence the critical
equations say that all three matrices commute.
On the ordered eigenvalue cover, put
$\delta=p_x-q_x$. The two spectral projectors are
\[
 e_p=\frac{A-q_xI}{\delta},\qquad
 e_q=\frac{p_xI-A}{\delta}.
\]
They are orthogonal idempotents of rank one, sum to the identity,
and commute with $B,C$. They define line bundles over every base
scheme on this cover. On each line, $B,C$ are scalar.
This constructs the inverse to the displayed quotient description,
including its automorphisms. Each invariant line is one of the
two eigenspaces. Recording that line orders the pair, which proves
the assertion about $\mathfrak E$.

The splitting also follows from the full extension complex.
For $R=\mathbb C[x,y,z]$, the Koszul resolution gives
$\operatorname{RHom}_R(\mathcal O_q,\mathcal O_p)$ as
$\bigwedge^\bullet\mathbb C^3$ with differential
\[
 d=v\wedge(-),\qquad
 v=(p_x-q_x)e_x+(p_y-q_y)e_y+(p_z-q_z)e_z.
\]
Contraction $h=\delta^{-1}\iota_{e_x^*}$ satisfies $dh+hd=1$.
Thus every cross extension group, including the cross Hom group,
vanishes. The opposite ordered pair has the same contraction with
its corresponding coordinate difference.
\end{proof}

\begin{lemma}[Transverse vanishing cycles and orientation]
\label{lem:c3-separated-vanishing}
There is an isomorphism
\[
 \left.\mathcal P_2\right|_{\mathfrak S}
 \simeq\mathbb Q_{\mathfrak S}[4],\qquad
 s^*\mathcal P_2\simeq
 \left.(\mathcal P_1\boxtimes\mathcal P_1)\right|_{\mathfrak D}.
\]
Choose its sign by giving the complex upper-triangular attracting
plane its Thom orientation. With this choice, exchanging the two
summands acts trivially on the rank-one transverse vanishing cycle.
\end{lemma}

\begin{proof}
Diagonalize $A$ on its ordered eigenvalue cover. The remaining
stabilizer is $\mathbb G_m^2$. The four transverse coordinates are
$b_{12},b_{21},c_{12},c_{21}$, and direct multiplication gives
\begin{equation}
 f_2=\delta(b_{12}c_{21}-b_{21}c_{12}).
 \label{eq:c3-separated-quadratic}
\end{equation}
In that order the Hessian is
\[
 \begin{pmatrix}
 0&0&0&\delta\\
 0&0&-\delta&0\\
 0&-\delta&0&0\\
 \delta&0&0&0
 \end{pmatrix},\qquad \det=\delta^4.
\]
The change of variables
\[
 u_1=b_{12},\quad u_2=c_{12},\quad
 v_1=\delta c_{21},\quad v_2=-\delta b_{21}
\]
gives $f_2=u_1v_1+u_2v_2$ over the entire ordered locus.
It uses $\delta^{-1}$ and no choice of its square root.
The Milnor fibre of this nondegenerate quadratic form retracts
onto $S^3$. Its reduced rational cohomology has rank one in degree
three. The convention $\phi_f^p=\phi_f[-1]$ therefore removes the four
normal dimensions from the ambient shift. The remaining stack has
dimension $6-2=4$, giving $\mathbb Q[4]$.
This is the quadratic stabilization calculation of
Brav--Bussi--Dupont--Joyce--Szendr\H{o}i,
Example~2.14 and Theorem~5.4 of \cite{BBDJSVanishing}.

The diagonal torus acts on each paired normal line and its dual
with opposite characters. Its determinant on the normal space is
one. Exchanging the eigenvalues exchanges both pairs
$b_{12},b_{21}$ and $c_{12},c_{21}$ and sends $\delta$ to $-\delta$.
Its determinant on the four original normal coordinates is $+1$.
The determinant of the displayed change to $(u_1,u_2,v_1,v_2)$ is
$-\delta^2$, unchanged under this exchange. The induced isometry
of $u_1v_1+u_2v_2$ consequently has determinant $+1$.
An orthogonal transformation acts on the middle vanishing cycle
by its determinant, as follows by its action on the oriented sphere.
The chosen generator therefore descends to $\mathfrak S$.
The same calculation trivializes the orientation local system.
\end{proof}

The Tate factor can also be recorded before forgetting mixed Hodge
structures. The affine quadric $u_1v_1+u_2v_2=1$ is isomorphic to
$SL_2(\mathbb C)$. Projection to its first column is an affine-line
bundle over $\mathbb C^2\setminus\{0\}$. The localization sequence
at the origin gives its third cohomology as $\mathbb Q(-2)$.
Accordingly, the untwisted mixed-Hodge version of the lemma reads
$s^*\mathcal P_2=(\mathcal P_1\boxtimes\mathcal P_1)(-2)$.
Multiplication in that coefficient theory must retain this Tate
factor, or specify a compensating dimension-dependent normalization.
The rational operations below use the stated forgetful convention.

The product can be defined on this open correspondence without
extending a class across the collision divisor. Set
\[
 H_D=H^*(\mathfrak D,\mathbb Q),\qquad
 H_S=H^*(\mathfrak S,\mathbb Q).
\]
The preceding shifts identify these with the corresponding critical
cohomology groups after a common regrading by four.
Use the complex Thom orientation in
Lemma~\ref{lem:c3-separated-vanishing}. The Hall operation $m:H_D\to H_S$
is induction along $b$ after the identification along $a$.
The ordered splitting map is $\Delta_{1,1}=s^*:H_S\to H_D$.
For two length-one classes $\alpha,\beta$, their local product is
$m((\alpha\boxtimes\beta)|_{\mathfrak D})$.

\begin{theorem}[Hall multiplication and ordered splitting]
\label{thm:c3-separated-hall}
On the separated critical locus, Hall induction is the finite-cover
trace $s_*$. The operations satisfy
\begin{equation}
 m\Delta_{1,1}=2\operatorname{id}_{H_S},\qquad
 \Delta_{1,1}m=\operatorname{id}_{H_D}+\tau^*.
 \label{eq:c3-separated-mackey}
\end{equation}
Both maps have rank one on the coefficient fibres
$\mathbb Q^2\rightleftarrows\mathbb Q$.
The splitting map is injective, and $m(1)=2$.
For the local system $\mathcal A=s_*\mathbb Q_{\mathfrak D}$ on
$\mathfrak S$, the relative pairing
\[
 \mathcal A\otimes\mathcal A\longrightarrow\mathbb Q_{\mathfrak S},
 \qquad (f,g)\longmapsto\operatorname{Tr}_s(fg)
\]
is perfect. Trace and pullback are adjoint for this pairing and
ordinary multiplication on $\mathbb Q_{\mathfrak S}$.
\end{theorem}

\begin{proof}
The smooth ambient extension stack consists of triples preserving
a line. Its map to $\mathfrak X_1\times\mathfrak X_1$ has relative
dimension $9-3-4=2$. The forgetful map is proper because the line
lies in the projective variety $\mathbb P^1$.
The potential on the extension stack is the sum of the two
diagonal potentials. This is the quiver induction correspondence
in \cite[\S1, equation~(1.1)]{KPS}.

In the ordered diagonal slice, preservation of the first eigenline
sets $b_{21}=c_{21}=0$. This is the plane $v_1=v_2=0$
in \eqref{eq:c3-separated-quadratic}. Hyperbolic restriction for the
cocharacter $\operatorname{diag}(t,1)$ integrates over this attracting
plane. Its complex Thom class generates
$H_c^4(\mathbb C^2,\mathbb Q)$ and has integral one.
The quadratic vanishing cycle has rank one, and its identification
with this Thom class is exactly the orientation fixed above.
Consequently the induction coefficient on each ordered critical
sheet is $+1$. The critical extension stack itself is $\mathfrak D$
by Lemma~\ref{lem:c3-separated-correspondence}. Proper pushforward
therefore sums the two sheet values. This proves $m=s_*$ and
identifies the operation with critical Hall induction in this
coefficient convention.

The fibre product of the cover with itself is the disjoint union
of the graphs of $1$ and $\tau$. Finite base change gives
$s^*s_*=1+\tau^*$. The trace of a pulled-back section is twice that
section, giving $s_*s^*=2$. On a trivializing open, the matrices are
\[
 m=\begin{pmatrix}1&1\end{pmatrix},\qquad
 \Delta_{1,1}=\begin{pmatrix}1\\1\end{pmatrix}.
\]
The relative pairing has matrix $\operatorname{diag}(1,1)$ in the
two sheet idempotents. This proves perfection and adjointness.
All these identities descend and hence hold on hypercohomology.
\end{proof}

\begin{proposition}[Equivariant cohomology and its trace form]
\label{prop:c3-separated-trace}
Let $u_i=c_1(L_i)$ have degree two. There is a degree-one class
$\eta$ such that
\[
 H_D=\mathbb Q[u_1,u_2]\otimes\bigwedge(\eta),\qquad
 H_S=\mathbb Q[c_1,c_2]\otimes\bigwedge(\eta),
 \quad c_1=u_1+u_2,\quad c_2=u_1u_2.
\]
Here $s^*$ is the displayed inclusion, $\tau^*\eta=\eta$, and
$m(f)=f+\tau^*f$. In particular,
\[
 m(u_1)=c_1,\qquad m(u_1u_2)=2c_2,\qquad
 \ker m=(u_1-u_2)H_S.
\]
As an $H_S$-module, $H_D$ is free with basis $1,d$, where
$d=u_1-u_2$. Its cohomological trace pairing has matrix
\[
 \begin{pmatrix}2&0\\0&2(c_1^2-4c_2)\end{pmatrix}.
\]
Thus this pairing is perfect after inverting $c_1^2-4c_2$.
It is not perfect over the displayed polynomial coefficient ring.
\end{proposition}

\begin{proof}
The coordinates $p_x-q_x,q_x,p_y,p_z,q_y,q_z$ identify $D$ with
$\mathbb C^*\times\mathbb C^5$. Its cohomology is the exterior
algebra on the winding class $\eta$ of the first coordinate.
Multiplication of that coordinate by $-1$ is homotopic to the
identity, so $\tau^*\eta=\eta$.
The Borel construction adds $\mathbb Q[u_1,u_2]$.
Rational descent through the finite group $S_2$ takes invariants
and has no higher group cohomology. The fundamental theorem of
symmetric polynomials gives the expression for $H_S$.
Every polynomial splits into its symmetric and antisymmetric
parts. The latter is divisible by $u_1-u_2$, with symmetric quotient.
This proves the module basis and kernel formula. Finally,
$d^2=c_1^2-4c_2$, $\operatorname{Tr}_s(d)=0$, and
$\operatorname{Tr}_s(1)=2$. The determinant is
$4(c_1^2-4c_2)$, which is a nonunit before localization.
\end{proof}

The last discriminant concerns equivariant Chern classes.
It differs from the spatial difference $p_x-q_x$, already inverted
in the definition of $D$. A perfect pairing on a finite-cover
coefficient sheaf can therefore induce a nonperfect pairing on
its equivariant cohomology. Integration over the nonproper base
and its $B\mathbb G_m$ stabilizers requires further support and
localization data before it yields a scalar Hall pairing.

In the point-brane interpretation, the Koszul contraction removes
open-string cross modes between separated branes. The four transverse
matrix coordinates form two quadratic pairs. The two terms of
$m(1)$ count the choices of a constituent line in their direct sum.
This interpretation uses the algebraic point-module sector and
does not identify it with renormalized holomorphic Chern--Simons observables.

If a global Hall coproduct restricts to $\Delta_{1,1}$, every
length-two primitive restricts to zero on $\mathfrak S$, by
\eqref{eq:c3-separated-mackey}. The collision locus requires new
data: at $p=q$, the cross Koszul differential is zero and its
cohomology dimensions are $1,3,3,1$.
The spectral projectors and the transverse quadratic reduction
then cease to apply. Extending multiplication and splitting across
collisions, comparing them with the chosen equivariant shuffle
coproduct, and constructing a continuous scalar Hopf pairing are
separate requirements for a geometric Hall double. The compact
Calabi--Yau comparison also requires orientation descent and a
map from the chosen field-theoretic observables.
